% ===== Closing Poem (七律) — to be set by the author by hand =====
\clearpage
\thispagestyle{empty}
\vspace*{2.2cm}
\begin{center}
{\small\color{gold} Coda}\\[0.3cm]
\textcolor{gold}{\rule{0.3\linewidth}{0.4pt}}\\[1.4cm]

% --- The 七律 is to be composed by the author (Wanhong), by hand. ---
% Discipline (consistent with the paper):
%   - a regulated verse (七律): eight lines, seven characters each,
%     with the two middle couplets in parallel (颔联 / 颈联);
%   - imagery-based, 含蓄, no closure spelled out (cf. Paper IX);
%   - no thesis terms ("generativity," "happiness," "symbolic," etc.);
%   - the anchor image of the rain beyond the window may return;
%   - the English below the poem is a plain prose gloss, not a verse
%     translation: it carries the sense for the non-Chinese reader while
%     the weight of the poem remains, untranslatable, in the Chinese
%     (which itself enacts §8: translation always leaves a remainder).
%
% Replace the placeholder block below with the finished poem.

{\cjk\large\color{rosedeep}
七律\\[0.7em]
\normalsize\color{rose}
闲窗相对坐忘言，　一霎回眸各已恬。\\[0.4em]
雨脚穿檐声未歇，　灯晕照座暖初添。\\[0.4em]
不须问处春常在，　无可名时意自圆。\\[0.4em]
同看廉纤千万缕，　此心如水落阶前。
}\\[1.6cm]

\begin{minipage}{0.78\linewidth}
\centering\itshape\footnotesize\color{inksoft}
By a quiet window, two sit facing each other, wordless; in a glance, each is already at ease. The rain runs along the eaves, its sound unceasing; the lamp's soft halo lights the seat, and a warmth is added. Where nothing need be asked, the spring is always there; in what cannot be named, the sense is round of itself. Together they watch the thousand fine threads of rain; this heart, like water, settles before the steps.
\end{minipage}

\vfill
\noindent\hfill\textcolor{gold}{\rule{0.2\linewidth}{0.4pt}}\hfill\null
\end{center}
\clearpage
